\chapter*{Введение}
\addcontentsline{toc}{chapter}{Введение}

Это мой личный конспект и собрание доказательств на естественном языке к книге Теренса Тао
\emph{Analysis I}. Конспект ведётся параллельно с формализацией той же книги на языке Lean~4
в этом же репозитории (директория \path{Analysis/}) и служит двум целям.

Во-первых, это способ разобраться в материале книги медленнее и внимательнее, чем позволяет
Lean: записать доказательство обычными словами — значит честно проверить, понимаю ли я, почему
утверждение верно. Тактика способна закрыть цель и без такого понимания.

Во-вторых, это мост между формальным и неформальным изложением. Каждое утверждение здесь
пронумеровано \emph{дословно} так же, как в оригинальной книге Тао, и снабжено ссылкой на
соответствующее имя в Lean-файле того же раздела — мелкая серая строка \texttt{Lean:~...} под
формулировкой. По ней всегда можно свериться, как одно и то же утверждение выглядит в прозе и
в машинно проверяемом коде.

\paragraph{Что здесь есть, а чего нет.}
Конспект не претендует на полноту и не заменяет книгу — это личные заметки читателя.
Не все определения, теоремы и упражнения из \emph{Analysis I} здесь появятся. В первую очередь
фиксируются доказательства, которые уже формализованы на Lean, и те места, где стоило
разобраться письменно, чтобы действительно понять рассуждение.

\paragraph{Как это читать.}
\begin{itemize}
  \item Нумерация (\emph{Аксиома 3.8}, \emph{Утверждение 2.1.4}, \dots) совпадает с
        нумерацией Тао — по ней легко найти утверждение в оригинальной книге.
  \item Строка \texttt{Lean:~...} под формулировкой — это имя леммы, теоремы или аксиомы в
        соответствующем файле \path{Analysis/Section_X_Y.lean}.
  \item Курсив используется по минимуму — только там, где есть настоящее смысловое ударение.
\end{itemize}
